\documentclass{ou}
\begin{document}
\thispagestyle{empty}
\begin{center}
\shortstack{\FGAVO}\vspace{4\baselineskip}
\textbf{РЕФЕРАТ}\\
\setstretch{2}
Информационные системы и~технологии\\
\setstretch{1}
Технологии создания и обработки табличных данных\vspace{2\baselineskip}\end{center}
\begin{flushleft}
\doublespacing
\begin{tabularx}{\textwidth}{@{} p{0.4\textwidth} l l @{}}
Руководитель & ст. преподаватель & И. Р. Нигматуллин\\
Студенты     &                   &                  \\
\multirow{-2.27}{\hsize}{\singlespacing Структура и организация табличных данных в~электронных таблицах}
             & ГФ25-01Б, 152536115 & С. В. Идрисова\\
             & ГФ25-01Б, 152540131 & Е. Кичко\\
             & ГФ25-01Б, 152539526 & К. Н. Аркадьева\\
             & ГФ25-01Б, 152536056 & В. А. Карасова\\
\multirow{-2.27}{\hsize}{\singlespacing Основные форматы хранения и~обмена табличной информацией}
             & ГФ25-01Б, 152537671 & А. Е. Корикова\\
             & ГФ25-01Б, 152536528 & Д. А. Васильева\\
             & ГФ25-01Б, 152540016 & И. Н. Морозова\\
\multirow{-2.27}{\hsize}{\singlespacing Анализ и визуализация данных средствами табличных приложений}
             & ГФ25-01Б, 152537684 & О. И. Бочанова\\
             & ГФ25-01Б, 152539986 & М. К. Гинтер\\
             & ГФ25-01Б, 152541471 & А. О. Безруких\\
             & ГФ25-01Б, 152536796 & А. В. Корюкова\\
\end{tabularx}
\end{flushleft}
\vspace*{\fill}
\begin{center}
\city
\end{center}
\newpage
\thispagestyle{empty}
\setstretch{1}
Продолжение титульного~листа реферата по~теме «Технологии создания и обработки табличных данных»
\begin{flushleft}
\doublespacing
\begin{tabularx}{\textwidth}{@{} p{0.4\textwidth} l l @{}}
Студенты     &                     &         \\
\multirow{-2.27}{\hsize}{\singlespacing Автоматизация расчетов и~использование формул и~функций}
             & ГФ25-01Б, 152541294 & Э. Р. Махмудов\\
             & ГФ25-01Б, 152539829 & В. Д. Бобкова\\
             & ГФ25-01Б, 152540753 & М. И. Бобылев\\
\multirow{-2.27}{\hsize}{\singlespacing Основы подготовки табличных данных для последующего анализа}
             & ГФ25-01Б, 152540016 & И. Н. Морозова\\
             & ГФ25-01Б, 152540765 & А. В. Гайворонская\\
             & ГФ25-01Б, 152537684 & О. И. Бочанова\\
             & ГФ25-01Б, 152539986 & М. К. Гинтер\\
             & ГФ25-01Б, 152538103 & С. Э. Манцаева\\
\multirow{-2.27}{\hsize}{\singlespacing Применение табличных инструментов в учебной и~исследовательской
работе}
             & ГФ25-01Б, 152540046 & А. Д. Балцевич\\
             & ГФ25-01Б, 152537659 & В. Е. Зайцева\\
             & ГФ25-01Б, 152536592 & А. Ю. Бирюкова\\
             & ГФ25-01Б, 152540765 & А. В. Гайворонская\\
            
\end{tabularx}
\end{flushleft}
\newpage
{
\sloppy
\tableofcontents
}
\newpage
\onehalfspacing
\section{Структура и организация табличных данных в~электронных таблицах}
Электронные таблицы -- это мощные средства для хранения, обработки и~анализа данных. Они широко используются в различных сферах деятельности: бизнесе, науке, образовании и других областях. Основным элементом электронных таблиц является таблица, которая позволяет структурировать данные, организовать их для дальнейшего анализа и автоматизации процессов. В данном реферате рассматриваются основные аспекты структуры и~организации табличных данных в электронных таблицах.
\subsection{Структура табличных данных}
Табличные данные представляют собой двумерную структуру, состоящую из строк и столбцов.

Строки -- горизонтальные объединения ячеек, каждая из которых содержит определённый набор данных или информации.

Столбцы -- вертикальные объединения ячеек, служат для~группировки однородных данных по определённым признакам.

Ячейка -- единичный элемент таблицы, который может содержать различные типы данных: числа, текст, даты, формулы и др. Адрес ячейки определяется её положением: например, «A1» (первая колонка, первая строка).

Диапазоны данных -- это совокупность нескольких ячеек, объединённых вместе (например, «A1:A10» или «B2:D5»), используется для групповой обработки данных.
\subsection{Организация данных в электронных таблицах}
Логическая структура:
\begin{enumerate}[leftmargin=0em, itemindent=3.75em, nosep, label=\arabic*)]
\item Данные структурированы по логике предметной области;
\item В заголовках столбцов задаётся описание признаков (например, «Имя», «Возраст», «Цена»);
\item Каждая строка обычно соответствует отдельному объекту или записи.
\end{enumerate}

Обеспечение целостности и правильности данных:
\begin{itemize}[label=-, leftmargin=0em, itemindent=3.35em, nosep]
\item Использование валидации данных (например, ограничение ввода чисел или дат);
\item Защита ячеек от случайных изменений;
\item Использование формул и функций для автоматической обработки.
\end{itemize}

В современных электронных таблицах реализована возможность преобразования диапазонов данных в структурированные таблицы (например, в Excel -- команда «Форматировать как таблицу»). Такая структура обеспечивает автоматические расширения, сортировку, фильтрацию и использование именованных диапазонов.

Основные принципы организации таблиц:
\begin{enumerate}[leftmargin=0em, itemindent=3.75em, nosep, label=\arabic*)]
\item Ясность и однозначность: структура таблицы должна быть понятна даже для человека, незнакомого с данными;
\item Однородность данных: в одном столбце хранятся одинаковые признаки;
\item Минимизация избыточности: избегать повторяющихся данных;
\item Использование формул и автоматизации: для сокращения ручной работы и повышения точности.
\end{enumerate}

Особенности организации данных в популярных электронных таблицах:
\begin{itemize}[label=-, leftmargin=0em, itemindent=3.35em, nosep]
\item Microsoft Excel: богатый набор инструментов для организации данных, сводные таблицы, функции проверки данных;
\item Google Таблицы: совместная работа, автоматическая синхронизация, встроенные скрипты;
\item LibreOffice Calc: открытое ПО с возможностями по~организации и~обработке таблиц.
\end{itemize}
\newpage
\section{Основные форматы хранения и обмена табличной информацией}
Электронная таблица -- это таблица данных, организованных в строки и столбцы. Очень часто используется для хранения финансовой информации благодаря возможности автоматически пересчитывать весь лист после изменения отдельной ячейки. Редактор электронных таблиц позволяет открывать, просматривать и редактировать самые популярные форматы файлов электронных таблиц. 
\subsection{Формат XLS}
XLS Расширение имени файла для электронных таблиц, созданных программой Microsoft Excel. XLS -- это проприетарный бинарный формат файлов, разработанный Microsoft для хранения данных электронных таблиц в программе Excel. Технически, XLS представляет собой сложную бинарную структуру, организованную согласно спецификации BIFF (Binary Interchange File Format). XLS-файлы хранят не только сами данные, но и метаданные о~форматировании, формулах, макросах, и других специфических элементах Excel. Бинарная природа формата делает его компактным, но одновременно и~менее прозрачным для сторонних программ. XLS считается устаревшим форматом, он сохраняет свою актуальность благодаря огромной базе исторических документов и системам, которые до сих пор используют этот формат для совместимости с устаревшим программным обеспечением. Многие современные программы продолжают поддерживать XLS для обеспечения обратной совместимости.

Ключевые преимущества формата XLS:
\begin{itemize}[label=-, leftmargin=0em, itemindent=3.35em, nosep]
\item Универсальная совместимость -- практически все программы для работы с электронными таблицами, включая устаревшие версии, поддерживают этот формат;
\item Меньший размер файла -- за счёт бинарной структуры XLS-файлы часто занимают меньше места, чем их XML-аналоги при небольших объёмах данных;
\item Полная интеграция с VBA -- макросы на Visual Basic for Applications полностью поддерживаются и выполняются без ограничений;
\item Стабильность формата -- спецификация не менялась годами, что гарантирует предсказуемое поведение при работе с файлами;
\item Поддержка устаревших систем -- многие корпоративные системы, ERP и банковское ПО оптимизированы именно под~формат XLS.
\end{itemize}

Существенные ограничения формата XLS:
\begin{enumerate}[leftmargin=0em, itemindent=3.75em, nosep, label=\arabic*)]
\item Лимит объёма данных -- невозможность работать с действительно большими массивами информации;
\item Ограниченная безопасность -- устаревшие механизмы шифрования делают защиту данных в XLS уязвимой;
\item Проблемы целостности -- бинарный формат более подвержен повреждениям, чем XML-форматы;
\item Новые типы данных, аналитические инструменты Excel часто несовместимы с XLS;
\item Ограничения по количеству цветов -- палитра из 56 цветов против миллионов в современных форматах.
\end{enumerate}

На практике выбор XLS обусловлен необходимостью взаимодействия с унаследованными системами или партнёрами, использующими устаревшее программное обеспечение. Для многих организаций переход на более современные форматы сопряжён с миграцией данных и~обновлением программного обеспечения, что требует значительных инвестиций.
\subsection{XLSX}
Формат XLSX стал стандартом для хранения и обработки электронных таблиц. Он используется в Microsoft Excel и поддерживается множеством других офисных приложений. В отличие от старого формата XLS, который применял бинарное представление данных, XLSX основан на открытом стандарте Open XML. Это делает его гибким, расширяемым и совместимым с различными программными платформами. XLSX Стандартное расширение для файлов электронных таблиц, созданных с помощью программы Microsoft Office Excel. 

Макрос в Excel -- это набор инструкций и команд, написанных на~языке программирования Visual Basic for Applications (VBA). Их используют для~автоматизации повторяющихся задач -- вместо того чтобы выполнять десяток повторяющихся действий, пользователь записывает одну команду и~затем запускает её, когда нужно совершить эти действия снова.

Преимущества формата XLSX:
\begin{enumerate}[leftmargin=0em, itemindent=3.75em, nosep, label=\arabic*)]
\item Формат XLSX стал стандартом для работы с электронными таблицами благодаря ряду ключевых преимуществ. Он сочетает в себе компактность, гибкость, безопасность и совместимость с различными системами;
\item Файл .xlsx представляет ZIP-архив, содержащий XML-документы. Это позволяет значительно уменьшать размер файлов по сравнению с устаревшим форматом .xls. Например, таблица на 10 000 строк с большим количеством формул может весить в 2-3 раза меньше, чем аналогичный файл в старом формате;
\item XLSX основан на Open XML, что делает его открытым и легко читаемым для сторонних программ. Любой разработчик может написать скрипт, который разберёт содержимое файла и извлечёт нужные данные без~использования Excel;
\item Формат XLSX поддерживается не только Microsoft Excel, но и другими офисными пакетами;
\item Форматирование ячеек -- В XLSX можно настраивать внешний вид данных, используя: изменение шрифтов, цветов и границ ячеек; условное форматирование (например, подсветку отрицательных значений), различные числовые форматы (валюта, проценты, даты);
\item Диаграммы и графики -- XLSX позволяет строить визуализации для наглядного представления данных: гистограммы, линейные и круговые диаграммы, диаграммы с областями, пузырьковые диаграммы, динамические графики, обновляемые при изменении данных;
\item Сводные таблицы -- этот инструмент позволяет группировать данные по категориям, автоматически рассчитывать суммы, средние значения, проценты, создавать интерактивные отчёты.
\end{enumerate}

Недостатки формата XLSX:
\begin{enumerate}[leftmargin=0em, itemindent=3.75em, nosep, label=\arabic*)]
\item Ограничения на объём данных -- Microsoft Excel накладывает ограничения на размер файлов: максимальное количество строк в одном листе -- 1 048 576; максимальное количество столбцов -- 16 384. Размер файла XLSX может достигать нескольких сотен мегабайт, что замедляет его открытие;
\item Проблемы с макросами -- хотя формат поддерживает макросы VBA, они могут работать не во всех программах. Файл XLSX с макросами должен быть сохранён в .xlsm, иначе VBA-код будет удалён;
\item Проблемы совместимости при конвертации -- если сохранить XLSX в другой формат (например, CSV или ODS), могут возникнуть ошибки: потеря форматирования ячеек и стилей; изменение кодировки, из-за чего русскоязычный текст может отображаться некорректно; пропадание формул при экспорте в CSV;
\item Риск повреждения файлов -- так как XLSX -- это ZIP-архив с~XML-файлами, его структура чувствительна к ошибкам. 
\end{enumerate}
\subsection{ODS}
ODS -- расширение имени файла для электронных таблиц, используемых пакетами офисных приложений OpenOffice и StarOffice, открытый стандарт для~электронных таблиц. 

Какие возможности поддерживаются при сохранении листа Excel в~формате электронной таблицы OpenDocument:
\begin{enumerate}[leftmargin=0em, itemindent=3.75em, nosep, label=\arabic*)]
\item Функция поддерживается и в формате Excel, и в формате электронной таблицы OpenDocument. Потери содержимого, форматирования не произойдет;
\item Функция частично поддерживается формате Excel, электронной таблицы OpenDocument, но возможно повреждение форматирования и снижение функциональности. Текст и данные не будут потеряны, но~форматирование и способ работы с текстом или графикой могут различаться;
\item Функция Excel не поддерживается в формате электронной таблицы OpenDocument. Если вы собираетесь сохранить лист Excel в формате электронной таблицы OpenDocument, не используйте эту функцию, иначе возможна потеря содержимого, форматирования и функциональности в~соответствующей части листа.
\end{enumerate}
\subsection{OTS}
Файл OTS -- это шаблон электронной таблицы, созданный программой Calc, входящей в состав Apache OpenOffice. Он содержит электронную таблицу, сохранённую в формате OASIS OpenDocument на основе XML. Пользователи Calc создают файлы OTS для копирования электронных таблиц с~теми же стилями и форматированием, что и в файлах ODS. OTS OpenDocument Spreadsheet Template -- формат текстовых файлов OpenDocument для шаблонов электронных таблиц. Шаблон OTS содержит настройки форматирования, стили и может использоваться для создания множества электронных таблиц со~схожим форматированием. 

Calc -- одна из нескольких программ, доступных в пакете OpenOffice. Она похожа на Excel, доступную в пакете Microsoft Office. В шаблонах OTS могут храниться различные типы электронных таблиц, например счета, бюджеты и календари. В Calc вы можете создавать собственные шаблоны с настраиваемыми макетами электронных таблиц, данными, шрифтами и стилями. Вы также можете скачать шаблоны с веб-сайта OpenOffice. Организация по развитию стандартов структурированной информации (OASIS) поддерживает стандарт OpenDocument.

Особенности формата в том, что он основан на XML. Состоит из~набора нескольких вложенных документов, каждый из которых хранит определённый аспект документа. Например, один вложенный документ содержит информацию о стиле, другой -- содержимое документа. По~умолчанию при открытии файла шаблона из системной оболочки открывается новый файл документа на основе шаблона.  

Файлы OTS используются для:
\begin{enumerate}[leftmargin=0em, itemindent=3.75em, nosep, label=\arabic*)]
\item Создания шаблонов с предопределёнными параметрами, связанными со стилями, шрифтом, данными, макетом электронной таблицы;
\item Создания нескольких файлов электронных таблиц ODS с заданным стилем и форматированием;
\item Определения настроек по умолчанию для различных типов документов, например, отгрузочных манифестов или бухгалтерских документов.
\end{enumerate}
\subsection{XLTX}
XLTX Excel Open XML Spreadsheet Template разработанный компанией Microsoft формат файлов на основе XML, сжатых по технологии ZIP. Предназначен для шаблонов электронных таблиц. Шаблон XLTX содержит настройки форматирования, стили и т.д. и может использоваться для создания множества электронных таблиц со схожим форматированием. 

Основные особенности XLTX:
\begin{enumerate}[leftmargin=0em, itemindent=3.75em, nosep, label=\arabic*)]
\item Шаблонность -- при открытии .xltx Excel не изменяет исходный файл, а создаёт новую копию (новую книгу .xlsx), основанную на нём. Это защищает шаблон от случайного изменения;
\item Хранение стилей и форматов -- цветовые схемы, шрифты, макеты листов, заголовки, логотипы, рамки, формулы и диаграммы можно заранее задать. Очень удобно для унификации отчётов, актов, финансовых таблиц;
\item Безопасность -- в .xltx нельзя хранить макросы (это отличие от~.xltm, где макросы разрешены). Благодаря этому шаблон безопасен для распространения внутри компаний;
\item Совместимость -- работает с Excel 2007+, LibreOffice Calc, Google Sheets (с частичными ограничениями). Хранится в виде ZIP-архива с~XML-файлами внутри.
\end{enumerate}
\subsection{XTML}
Формат XLTM -- это формат шаблона Microsoft Excel с поддержкой макросов. Он предназначен для создания новых файлов Excel, в которых заранее определены структура, стили, формулы и автоматизация через макросы VBA (Visual Basic for Applications). Макросы позволяют выполнять рутинные задачи в один клик. Можно задавать собственные функции, кнопки, меню, обработчики событий (например, при открытии или изменении ячеек).

XLTM и его преимущества:
\begin{itemize}[label=-, leftmargin=0em, itemindent=3.35em, nosep]
\item Шаблон Excel, который можно использовать многократно;
\item Позволяет хранить макросы VBA для автоматизации;
\item Используется для типовых, но сложных таблиц, где требуется не только структура, но и логика обработки;
\item Удобен для корпоративных отчётов, финансов, аналитики.
\end{itemize}
\subsection{CSV: особенности, применение и причины создания}
Файл CSV -- это текстовый файл, в котором каждая строка соответствует одной строке таблицы, а значения в ней разделяются определённым символом -- чаще всего запятой. Однако разделителем может быть и другой символ, например точка с запятой, табуляция или пробел.

Основные особенности CSV:
\begin{enumerate}[leftmargin=0em, itemindent=3.75em, nosep, label=\arabic*)]
\item Простота и прозрачность. CSV-файлы -- это обычный текст, который можно прочитать любым редактором;
\item Минимум ограничений. Нет строгих требований к длине строки, количеству столбцов или типу данных;
\item Гибкость разделителей. Символ разделения может быть любым, главное -- чтобы он не встречался внутри самих данных;
\item Отсутствие метаданных. CSV хранит только «сырые» значения без~информации о типах данных, форматировании или формулах;
\item Совместимость. CSV поддерживается большинством офисных и~аналитических приложений, включая Excel, LibreOffice, Google Sheets, R, Python, SQL.
\end{enumerate}

Преимущества формата:
\begin{enumerate}[leftmargin=0em, itemindent=3.75em, nosep, label=\arabic*)]
\item CSV не требует сложных алгоритмов для~чтения или записи. Программисту достаточно разбить строку по~символу-разделителю, обработать полученные значения. Благодаря этому CSV можно использовать в любых языках программирования без дополнительных библиотек;
\item Независимость от платформы -- так как это текстовый формат, он одинаково хорошо работает на всех операционных системах -- Windows, macOS, Linux. Это делает CSV идеальным средством передачи данных между разными средами;
\item Совместимость с офисными пакетами -- Microsoft Excel и другие таблицы позволяют открывать и сохранять файлы CSV без дополнительных инструментов. Таким образом, данные, созданные в одном приложении, можно легко открыть в другом;
\item Минимальный размер и высокая скорость обработки -- по сравнению с бинарными форматами, CSV занимает немного места и быстро читается. Это особенно важно для обмена большими наборами данных по сети;
Прозрачность и долговечность -- CSV-файлы не зависят от конкретной версии программы. Их можно открыть и через десятки лет, не боясь, что формат станет устаревшим.
\end{enumerate}

Ограничения CSV:
\begin{enumerate}[leftmargin=0em, itemindent=3.75em, nosep, label=\arabic*)]
\item Отсутствие строгого стандарта. Разные программы могут по-разному трактовать разделители, кавычки и кодировки. Это создаёт проблемы при~импорте данных;
\item Нет поддержки типов данных. Все значения хранятся как текст. Информация о числовом или временном типе теряется;
\item Нет поддержки иерархий и вложенных структур. CSV подходит только для плоских таблиц, проблемы с кодировками. Файлы, созданные в Windows (например, в CP1251), не всегда корректно открываются в Linux, где ожидается UTF-8;
\item CSV не хранит информацию о форматировании, ширине колонок, формулах.
\end{enumerate}

Формат CSV -- одно из тех решений, простота которых оказалась ключом к долговечности. Созданный как средство обмена табличными данными между несовместимыми системами, он пережил десятки технологических эпох. CSV демонстрирует, что минимализм и открытость могут быть не менее ценными, чем сложные стандарты. Он не обладает сложными структурами, не поддерживает типизацию, не хранит метаданные, но именно в этом его сила: любой человек или программа могут создать и прочитать CSV без~предварительных соглашений. Сегодня CSV используется в бизнесе, науке, образовании и государственных проектах. Его роль как «универсального формата» остаётся важной, а способность сохранять совместимость между системами делает его одним из наиболее устойчивых инструментов цифрового мира.
\subsection{PDF в работе с табличными данными}
С развитием цифровых документов возникла необходимость в формате, способном сохранять внешний вид документа на любых устройствах и в~любых операционных системах. Таким форматом стал PDF (Portable Document Format). PDF был создан как инструмент для точного воспроизведения документов, включающих текст, изображения, графику и таблицы. Его основная идея -- сохранить документ в том виде, в котором он был создан, независимо от~платформы, программного обеспечения или оборудования. Это качество сделало формат PDF не только стандартом для обмена официальными документами, но и широко используемым средством представления табличных данных -- отчётов, финансовых сводок, статистики, спецификаций и прочих структурированных сведений.

PDF представляет странично-ориентированный формат, основанный на~модели описания содержимого страницы. В отличие от текстовых форматов, где информация хранится в виде последовательности символов, PDF сохраняет координаты и свойства каждого элемента.
\subsection{Представление таблиц в формате PDF. Визуальная модель таблицы} 
Таблица в PDF -- это совокупность: линий, определяющих границы ячеек; текстовых объектов, расположенных в определённых координатах; графических элементов для форматирования (цвет фона, шрифт, выравнивание). PDF не содержит специальной конструкции «table» (в отличие от HTML). Каждая ячейка создаётся вручную при генерации документа программами, такими как Excel, Word, ReportLab, LaTeX и другими.

Причины использования PDF для табличных данных:
\begin{enumerate}[leftmargin=0em, itemindent=3.75em, nosep, label=\arabic*)]
\item Гарантированное визуальное соответствие -- PDF сохраняет документ в точности таким, каким он был создан, включая шрифты, цвета, размеры ячеек и форматирование;
\item Кроссплатформенность -- PDF одинаково отображается на Windows, macOS, Linux, Android и iOS. Получатель документа видит таблицу именно так, как её оформил автор, независимо от программ и устройств;
\item Универсальность и распространённость -- почти каждое современное устройство и браузер поддерживают PDF «из коробки». Это делает его идеальным форматом для обмена готовыми документами, не требующими редактирования;
\item Безопасность и защита данных -- PDF поддерживает шифрование, цифровые подписи и права доступа. Это позволяет защищать табличные документы от несанкционированного редактирования или копирования;
\item Поддержка сложной вёрстки -- PDF может содержать не только таблицы, но и сопроводительные тексты, графики, диаграммы, изображения. Это делает его удобным для комплексных отчётов, где табличные данные -- лишь часть визуальной структуры.
\end{enumerate}

Недостатки:
\begin{enumerate}[leftmargin=0em, itemindent=3.75em, nosep, label=\arabic*)]
\item Отсутствие редактируемой структуры. Нельзя напрямую работать с~данными как в Excel;
\item Сложность извлечения информации. Таблицы в PDF плохо поддаются автоматическому анализу;
\item Больший размер файла. Из-за графической природы PDF может быть объёмным;
\item Ограниченная машиночитаемость. Без тегов или дополнительных данных программа не понимает, где строки и столбцы.
\end{enumerate}

Формат PDF, изначально созданный для точного воспроизведения документов, стал неотъемлемой частью работы с табличной информацией. Хотя он не предназначен для хранения данных в структурированном виде, его способности сохранять оформление и кроссплатформенность сделали его стандартом в сфере представления таблиц. PDF идеально подходит для~финальной стадии жизненного цикла данных -- публикации и~распространения.

Таким образом, роль PDF в контексте таблиц -- не быть инструментом анализа, а стать гарантией неизменности и достоверности представленных данных. Его использование остаётся важным элементом цифрового документооборота, особенно там, где визуальная форма таблицы столь же значима, как и её содержимое.
\newpage
\section{Анализ и визуализация данных средствами табличных приложений}
Каждый день появляется море информации -- от компаний, магазинов, больниц, школ и просто людей. Чтобы не утонуть в этих данных, нужны простые и мощные инструменты.

Табличные приложения -- это программы вроде Microsoft Excel, Google Таблиц и OpenOffice Calc. Они предоставляют пользователям широкий набор инструментов для систематизации, расчётов, построения графиков и~представления данных в наглядной форме.

Анализ данных -- это процесс систематического изучения информации с целью выявления закономерностей, тенденций и взаимосвязей между различными показателями. Основная цель анализа данных -- получение новых знаний, которые могут быть использованы для принятия обоснованных решений.

В рамках табличных приложений анализ данных реализуется посредством формул, функций, фильтров, сводных таблиц и инструментов статистической обработки.
\subsection{Что умеют табличные приложения}
Табличные приложения (электронные таблицы) умеют создавать, редактировать и анализировать данные в виде таблицы. Это компьютерный эквивалент обычной таблицы, в ячейках которой записаны данные различных типов: текст, даты, формулы, числа. Каждая ячейка имеет уникальное адресное обозначение, позволяющее обращаться к ней и использовать содержащиеся в~ней значения в других ячейках.

Функции:
\begin{enumerate}[leftmargin=0em, itemindent=3.75em, nosep, label=\arabic*)]
\item Автоматизация расчётов с помощью формул, функций. При~изменении исходных данных все результаты автоматически пересчитываются и заносятся в таблицу;
\item Форматирование -- изменение цвета ячеек, шрифтов и их размеров, разные стили: от рамок до фоновой заливки;
\item Сортировка и фильтры -- упорядочивание информации в соответствии с критериями: по возрастанию или по убыванию. Фильтры скрывают лишние ячейки и оставляют только нужные данные;
\item Построение графиков и диаграмм -- создание столбчатых, круговых, линейных и других графиков, настройка их внешнего вида, добавление подписей;
\item Импорт и экспорт данных -- возможность загружать данные различных типов со сторонних ресурсов;
\item Интеграция с другими программами -- возможность вставлять текст, рисунки, таблицы, подготовленные в других приложениях. 
\end{enumerate}
\subsection{Сравнение Microsoft Excel, Google Таблиц и OpenOffice Calc}

Microsoft Excel, Google Таблицы и OpenOffice Calc -- программы и~сервисы для работы с электронными таблицами, но имеют разные особенности. Выбор зависит от задач: Excel ориентирован на локальную работу, Google Таблицы -- на облачную, а OpenOffice Calc -- на работу с~таблицами с открытым исходным кодом.  

Microsoft Excel --- программа для работы с электронными таблицами, часть пакета Microsoft Office. Некоторые возможности:  
\begin{enumerate}[leftmargin=0em, itemindent=3.75em, nosep, label=\arabic*)]
\item Создание и редактирование таблиц, добавление или удаление строк и~столбцов;
\item Форматирование таблиц: выделение ячеек, применение стилей и~цветов, изменение шрифта;
\item Формулы и функции для автоматизации вычислений, например, суммирование значений, вычисление среднего значения;
\item Инструменты для анализа данных: создание графиков и диаграмм, фильтры и сортировка;
\item Импорт и экспорт данных: можно импортировать данные из других программ или файлов формата CSV или XML, а также экспортировать таблицы и графики.
\end{enumerate}

Google Таблиц -- онлайн-сервис для создания и редактирования электронных таблиц, входит в состав Google Workspace. Некоторые особенности:  
\begin{itemize}[label=-, leftmargin=0em, itemindent=3.35em, nosep]
\item Совместная работа: пользователи могут одновременно редактировать документ, видеть изменения в реальном времени и оставлять комментарии;
\item Функции и формулы: сервис поддерживает более 500 функций, включая математические, статистические, текстовые и логические;
\item Визуализация данных: можно создавать диаграммы и графики прямо в~документе;
\item Интеграция с другими сервисами Google: сервис интегрирован с Google Документы, Google Формы, Google Диск и другими.
\end{itemize}

OpenOffice Calc -- табличный процессор, входящий в состав Apache OpenOffice и OpenOffice.org. Некоторые возможности: 
\begin{enumerate}[leftmargin=0em, itemindent=3.75em, nosep, label=\arabic*)]
\item Пошаговый ввод формул в ячейки электронных таблиц с помощью Мастера, который демонстрирует описания каждого параметра и конечный результат на любом этапе ввода;
\item Условное форматирование и стили ячеек позволяют упорядочить готовые данные;
\item Более двух десятков форматов импорта и экспорта файлов, включая функции импорта текста;
\item Поддерживаются связи между разными электронными таблицами и~совместное редактирование данных (начиная с версии OpenOffice.org 3.0).
\end{enumerate}
\subsection{Средства визуализации данных в табличных приложениях}
Наиболее распространенным средством визуализации в табличных приложениях являются графики и диаграммы. Гистограммы и столбчатые диаграммы полезны для сравнения количественных данных по различным категориям. Круговые диаграммы показывают долю каждой категории относительно общего объема. Линейные графики демонстрируют изменения данных во времени. Диаграммы рассеяния используются для выявления взаимосвязей между двумя переменными.

Такие инструменты позволяют преобразовывать сложные числовые данные в понятные и легкодоступные визуальные образы.

Некоторые табличные приложения поддерживают интерактивные элементы, увеличивающие уровень взаимодействия с~визуализацией:
\begin{enumerate}[leftmargin=0em, itemindent=3.75em, nosep, label=\arabic*)]
\item Фильтры и сортировка дают возможность динамического отбора данных для просмотра нужной информации;
\item Сводные таблицы и графики позволяют соединять и сравнивать большие объёмы данных;
\item Интерактивные карты и географические диаграммы -- визуализация пространственного расположения объектов и распределение данных по~территории.
\end{enumerate}

Эти элементы помогают пользователям взаимодействовать с данными и~получать детальную информацию по мере необходимости.
\subsection{Дополнительные инструменты и интеграции}
VBA (Visual Basic for Applications) -- встроенный язык программирования в Excel, позволяющий разрабатывать собственные макросы и модули для~автоматизации задач и формирования нестандартных визуализаций.

Библиотеки Python (Matplotlib, Seaborn) -- интеграция табличных приложений с языками программирования позволяет расширить диапазон возможных визуализаций и повысить точность представления данных.
\newpage
\section{Автоматизация расчетов и использование формул и функций}
Автоматизация в Excel -- это практический подход к организации работы с данными, который позволяет заменить ручные повторяющиеся операции на системные решения. Основная цель -- не просто ускорить процессы, но~и~повысить их точность, снизить вероятность ошибок и создать устойчивую среду для работы с информацией.
\subsection{Инструменты автоматизации}
В Excel доступны несколько инструментов для автоматизации. Формулы и функции служат базовым уровнем -- они позволяют автоматизировать вычисления, проверки и преобразования данных. Например, с помощью функций ВПР или СУММЕСЛИ можно создавать сложные расчетные системы, которые мгновенно обновляются при изменении исходных данных.

Power Query представляет более продвинутый инструмент -- он специализируется на сборе, очистке и объединении данных из разных источников. Это решение особенно ценно, когда нужно регулярно обрабатывать большие объемы информации из нескольких файлов или баз данных.

Для сложных многошаговых процессов предназначены макросы и~VBA. Они позволяют записывать последовательности действий, создавать пользовательские интерфейсы и автоматизировать задачи, которые невозможно решить стандартными средствами Excel.
\subsection{Практическая реализация}
Процесс внедрения автоматизации начинается с анализа. Нужно выявить задачи, которые повторяются регулярно и отнимают значительное время. Это может быть ежедневное составление отчетов, ручная проверка данных или~сложные расчеты, выполняемые по одинаковой схеме.

На этапе разработки подбираются подходящие инструменты. Простые задачи часто решаются формулами, для работы с данными из нескольких источников выбирается Power Query, а сложные процессы автоматизируются через макросы. Важно создать прототип решения и проверить его на реальных данных.

Тестирование -- критически важный этап. Автоматизированная система должна корректно работать не только в идеальных условиях, но~и~при~возникновении ошибок. Проверяются различные сценарии: некорректный ввод, отсутствие данных, граничные значения. Результаты автоматических расчетов сравниваются с ручными методами для обеспечения точности.

Обучение пользователей завершает процесс внедрения. Даже самая совершенная система не будет эффективной, если сотрудники не понимают, как с ней работать. Необходимо подготовить инструкции, провести обучение и организовать поддержку на начальном этапе.
\subsection{Области применения}
В финансовой сфере автоматизация Excel используется для составления отчетов, анализа инвестиций и бюджетирования. В операционной деятельности -- для управления запасами, контроля производственных показателей и~оптимизации цепочек поставок. В аналитике -- для создания дашбордов, отслеживания ключевых метрик и выявления тенденций.
\newpage
\section{Основы подготовки табличных данных для последующего анализа}
Подготовка табличных данных для последующего анализа включает очистку, форматирование и использование инструментов для организации информации.

В Excel:
\begin{enumerate}[leftmargin=0em, itemindent=3.75em, nosep, label=\arabic*)]
\item Удалить ненужные столбцы и строки -- лишние элементы мешают сосредоточиться на главном и замедляют обработку;
\item Сделать заголовки столбцов уникальными и описать их содержимое -- это позволяет быстро найти нужную информацию;
\item Проверить консистентность форматов данных в столбцах -- использовать единый формат для каждого вида данных. Например, для~дат -- все значения записаны в формате ДД/ММ/ГГГГ;
\item Использовать условное форматирование для создания подсветки ключевых элементов -- это позволяет быстро визуально идентифицировать важные записи;
\item Применить сортировку и фильтрацию -- сортировка помогает организовать данные, фильтрация позволяет сосредоточиться на нужных данных. Например, для сортировки нужно выделить нужный диапазон и~выбрать «Сортировка» в меню «Данные».
\end{enumerate}

В Google Sheets:
\begin{enumerate}[leftmargin=0em, itemindent=3.75em, nosep, label=\arabic*)]
\item Организовать данные структурированно -- каждый столбец должен представлять отдельную переменную, а каждая строка -- отдельное наблюдение или период. Такой подход упростит последующий анализ и визуализацию;
\item Использовать фильтры -- они позволяют находить в таблице нужные данные, показывать или скрывать определённые строки. Например, можно создать фильтр «Основание группировки», с его помощью можно группировать строки по выбранному полю;
\item Создать сводные таблицы -- они позволяют систематизировать данные, выявлять закономерности и упорядочивать информацию. Чтобы создать сводную таблицу, нужно перейти во вкладку «Данные», выбрать «Сводная таблица», настроить диапазоны данных и выбрать поля для анализа;
\item Сводная таблица -- это эффективный инструмент для вычисления, анализа данных, который упрощает поиск сравнений, закономерностей и~тенденций. Сводные таблицы работают немного по-разному в зависимости от платформы, используемой для запуска Excel.
\end{enumerate}

Советы и рекомендации по форматированию данных:
\begin{enumerate}[leftmargin=0em, itemindent=3.75em, nosep, label=\arabic*)]
\item Используйте чистые табличные данные для достижения наилучших результатов;
\item Убедитесь, что все столбцы имеют заголовки с одной строкой уникальных, непустых меток для каждого столбца;
\item Отформатируйте данные как таблицу Excel (выберите в любом месте данных, а затем на ленте выберите Вставить > таблицу);
\item Если у вас есть сложные или вложенные данные, используйте Power Query для их преобразования (например, для отмены сворачивания данных), чтобы они были упорядочены по столбцам с одной строкой заголовка.
\end{enumerate}

Сводные таблицы в Excel подходят для анализа данных. Они берут информацию из обычных таблиц, разбивают ее на блоки, выполняют необходимые вычисления, а затем представляют полученный результат наглядно. Кроме того, все показатели этого отчета могут быть настроены в~соответствии с требованиями пользователя.

Какие преимущества? Создав сводную таблицу, можно вносить в~нее изменения и дополнения, просто выделяя различные ячейки или~корректируя формулы в окне редактора. Это гораздо проще, чем вручную обновлять каждую ячейку по отдельности, что может занять много времени в зависимости от~объема данных.

Вот какие возможности они предоставляют:
\begin{itemize}[label=-, leftmargin=0em, itemindent=3.35em, nosep]
\item Структурирование данных;
\item Группировка и категоризация;
\item Фильтрация и сортировка;
\item Изменение структуры;
\item Расчет итогов;
\item Разворачивание данных;
\item Публикация и представление.
\end{itemize}

При работе со сводными таблицами/реестрами необходимо уделять внимание подготовке исходных данных. Вот несколько важных требований, которые следует учитывать:
\begin{enumerate}[leftmargin=0em, itemindent=3.75em, nosep, label=\arabic*)]
\item Каждый столбец в исходной таблице должен иметь соответствующий заголовок, чтобы обеспечить ясность и понимание данных;
\item Все значения в одной колонке должны быть введены в едином формате. Например, если у вас есть столбец «День поставки» то все данные должны быть записаны в формате даты. Аналогично, колонка «Поставщик» должна содержать только названия компаний или коды снабженцев, но не~смешанные форматы;
\item Исходная таблица не должна содержать полностью пустых строк или~столбцов, так как это может вызвать проблемы при анализе и создании сводной;
\item В ячейки следует вводить атомарные значения, то есть такие, которые нельзя разделить на более мелкие части и разнести по разным столбцам. Например, адрес следует разбивать на отдельные столбцы, такие как «Город» «Улица» и «Дом» чтобы обеспечить эффективную работу со сводными данными;
\item Избегайте создания таблиц с неправильной структурой. Старайтесь организовывать данные так, чтобы они отражали естественную схему информации и соответствовали целям анализа. 
\end{enumerate}

Исходные данные для создания сводных таблиц:
\begin{enumerate}[leftmargin=0em, itemindent=3.75em, nosep, label=\arabic*)]
\item Убедитесь, что информацию можно представить в виде таблицы, где строки представляют собой записи или наблюдения, а столбцы -- разные их атрибуты/характеристики;
\item Выделите все используемые ячейки данных и перейдите на вкладку «Главная» в меню программы, затем выберите «Форматировать как таблицу»;
\item Добавьте уникальные и информативные заголовки в каждый столбец. Они будут использоваться в сводном отчете в качестве имен полей;
\item Предварительно убедитесь, что исходная таблица не содержит пустых строк или столбцов, а также избавьтесь от промежуточных итогов, которые могут исказить результаты;
\item Для упрощения работы и обращения к первичным данным дайте таблице уникальное имя, которое можно ввести в поле «Имя» в верхнем правом углу.
\end{enumerate}
\subsection{Упорядочивание полей}
Можно изменить порядок отображения полей в сводной таблице тремя различными способами:
\begin{itemize}[label=-, leftmargin=0em, itemindent=3.35em, nosep]
\item Перетащите поле с помощью мыши между четырьмя областями размещения. Либо щелкните и удерживайте его имя в разделе «Поле» и~переместите его в нужную область в разделе «Макет»;
\item Правой кнопкой мыши кликните на имени поля в разделе «Поле» и~выберите нужную область для перемещения;
\item Нажмите на поле в разделе «Макет», чтобы выбрать его. Это позволит моментально просмотреть доступные параметры.
\end{itemize}

Все внесенные исправления применяются немедленно. Если случайно сделали что-то неправильно, не забудьте о комбинации клавиш CTRL+Z, которая позволяет отменить последние изменения (если, конечно, вы не~сохранили их, нажав соответствующую клавишу).

После создания сводной таблицы важно помнить, что если в исходный реестр были добавлены новые записи (строки, столбцы), то эти данные автоматически не появятся в ней. 

\subsection{Функции для значений}
В Microsoft Excel по умолчанию используется функция «Сумма» для~числовых данных, размещаемых в области «Значения». Если в эту область помещаются нечисловые данные (текст, даты или логические значения) или ячейки с пустыми данными, то автоматически применяется функция «Количество».

Выбор другого метода вычислений -- это возможность, которая у вас всегда есть. Имя поля можно легко изменить на то, которое будет понятно и~информативно. Читаемость и содержательность -- одни из важных аспектов имени, отображаемого в таблице.

Еще одной полезной возможностью является представление данных разными способами, например, в процентах или в виде ранжированных значений от наименьшего к наибольшему и наоборот. Эту функцию называют «Дополнительные вычисления». Для ее использования, как описано выше, просто перейдите на вкладку «Параметры …».

Совет: Функция «Дополнительные вычисления» особенно полезна, когда нужно добавить одно и то же поле более одного раза и отобразить, например, общий показатель продаж и их сумму продаж в процентах от общего объема одновременно. Используя эту функцию, вы экономите много времени, которое раньше приходилось бы тратить на создание сложных формул.

\subsection{Примеры использования сводных таблиц}
Современный мир данных нашел надежного спутника в лице Excel. Великолепные сводные таблицы стали верным проводником в лабиринте больших объемов информации. С их помощью можно умело анализировать данные, находя новые ответы на привычные вопросы.

Представьте, что сводные таблицы в Excel применяются для анализа продаж товаров в разные регионы. Вы можете узнать, какие товары приносят максимальную выручку, какие категории завоевали наибольшую популярность, и что именно привлекает внимание клиентов.

Сводные таблицы в Microsoft Excel помогут в анализе данных о~клиентах. Группируя их по разным параметрам, таким как возраст, пол или~местоположение, можно получить четкое представление о целевой аудитории. Это позволит лучше понимать потребности клиентов, разрабатывать более точные маркетинговые стратегии.

Но это еще не всё. Сводные таблицы в Excel также могут быть применены для анализа данных о производстве. С их помощью легко подытожить информацию о производственных процессах, группировать ее по периодам или этапам, и таким образом определить наиболее эффективные или, наоборот, требующие улучшений процессы.

Excel и его сводные таблицы становятся незаменимым инструментом изучения данных и помогают бизнес-аналитикам и компаниям принимать важные решения и оптимизировать процессы, всегда предоставляя ценную информацию в нужный момент.

Сводные таблицы в Excel представляют собой мощный инструмент анализа данных в Excel. Они решают множество задач, включая создание сводок, исследование, изучение и визуализацию информации. Дополняясь диаграммами, они позволяют не только обобщить показатели, но и наглядно отобразить тренды и сравнения. Сводные таблицы в Excel и графики помогают организации принимать обоснованные решения, опираясь на критически важные данные.

Иногда бывает необходимо преобразовать сводную таблицу в обычную форму. Для этого достаточно выполнить несколько простых шагов. Выберите любую ячейку в первом реестре, затем перейдите в меню «Работа со сводными таблицами», выберите «Конструктор» и в подменю «Макет отчета» кликните опцию «Показать в табличной форме». Теперь сводная таблица стала удобным обычным реестром.

Сводные таблицы в Microsoft Excel -- незаменимый инструмент анализа данных, облегчающий обобщение информации и помогающий в принятии решений на ее основе. С их помощью можно быстро и наглядно оценить тенденции и паттерны в больших наборах сведений.
\newpage
\section{Применение табличных инструментов в учебной и исследовательской работе}
Таблицы -- универсальный инструмент структурирования данных, который широко применяется как в образовательном процессе, так и в~научных исследованиях.

Основные функции таблиц:
\begin{itemize}[label=-, leftmargin=0em, itemindent=3.35em, nosep]
\item Систематизация информации -- упорядочивание разнородных данных по строкам и столбцам;
\item Сравнение параметров -- наглядное сопоставление характеристик объектов;
\item Выявление закономерностей -- обнаружение тенденций и корреляций;
\item Оптимизация объёма -- сжатое представление больших массивов данных;
\item Повышение наглядности -- визуализация количественных показателей.
\end{itemize}
\subsection{Сферы применения}
В учебной работе:
\begin{itemize}[label=-, leftmargin=0em, itemindent=3.35em, nosep]
\item Оформление результатов лабораторных работ;
\item Составление сводных таблиц по темам;
\item Систематизация теоретического материала;
\item Подготовка отчётов по практикам;
\item Создание конспектов и шпаргалок;
\item Построение графиков на основе табличных данных.
\end{itemize}

В учебной работе таблицы используют для организации и представления информации. Они помогают улучшить читаемость и понимание материала, так как предоставляют данные в компактной и структурированной форме. Таблицы также позволяют легко сопоставить и проанализировать информацию, выделить ключевые моменты.

В исследовательской деятельности:
\begin{itemize}[label=-, leftmargin=0em, itemindent=3.35em, nosep]
\item Фиксация первичных данных эксперимента;
\item Статистическая обработка результатов;
\item Сравнительный анализ методик;
\item Представление итоговых показателей;
\item Формирование приложений к научным работам;
\item Подготовка материалов для публикаций.
\end{itemize}

Таблицы в исследовательских работах служат для систематизации данных, наглядного сравнения показателей и представления количественных результатов, что позволяет выявить закономерности и подкрепить выводы эмпирическими доказательствами. Они экономят текстовое пространство, повышают доступность информации и обеспечивают воспроизводимость исследования. Кроме того, таблицы помогают соответствовать научным стандартам оформления и делают аргументацию более строгой и прозрачной.

Преимущества таблиц:
\begin{enumerate}[leftmargin=0em, itemindent=3.75em, nosep, label=\arabic*)]
\item Ясность и компактность. Таблицы позволяют представить большой объем информации в лаконичной форме. Без длинных абзацев текста -- таблицы делают информацию легкоусвояемой;
\item Сравнение и анализ. Столбцы и строки таблицы -- удобство для~анализа. Вы можете сравнить данные, выделить тренды, и выявить важные паттерны, всего лишь глядя на таблицу;
\item Упорядочивание данных. Таблица -- это шкафчик для ваших данных. Каждая ячейка хранит четко упорядоченные данные, не позволяя им путаться и разбредаться;
\item Визуальное воздействие. Подчеркните свою точку зрения, используя структурированные таблицы. Они могут быть мощным инструментом в~убеждении аудитории и делают вашу курсовую работу более качественной;
\item Экономия места. Компактное размещение большого объёма информации;
\item Удобство восприятия. Чёткая структура облегчает понимание;
\item Быстрый поиск. Возможность оперативно находить нужные данные;
\item Наглядность. Визуальное выделение ключевых показателей;
\item Универсальность. Подходит для любых типов данных (числовых, текстовых, символьных).
\end{enumerate}

Типичные ошибки при работе с таблицами:
\begin{itemize}[label=-, leftmargin=0em, itemindent=3.35em, nosep]
\item Отсутствие нумерации и заголовков;
\item Избыточное количество линий;
\item Нечитаемые сокращения;
\item Несогласованность единиц измерения;
\item Дублирование информации в тексте и таблице;
\item Нарушение логической структуры.
\end{itemize}

Рекомендации по эффективному использованию:
\begin{enumerate}[leftmargin=0em, itemindent=3.75em, nosep, label=\arabic*)]
\item Включайте в таблицу только релевантные данные;
\item Соблюдайте единообразие форматирования;
\item Сопровождайте таблицу аналитическим комментарием;
\item Используйте цветовое кодирование для выделения ключевых значений;
\item Проверяйте корректность расчётов перед публикацией;
\item Учитывайте требования ГОСТ или методических рекомендаций учебного заведения.
\end{enumerate}
\newpage
{
\sloppy
\nocite{*}
\cleardoublepage
\addcontentsline{toc}{section}{Список использованных источников}
\bibliographystyle{plainurl}
\bibliography{report5}
}
\end{document}
